% BEGIN COMPLETE CONSTRUCTION: this block may be placed in an appendix.
\section{Complete construction for budget hardness}
\label{hard:construction}

\paragraph{Source problem.}
An instance of two-budget bipartite partial vertex cover consists of a
simple bipartite graph \(H=(X,Y,E_H)\), nonnegative integer budgets \(a,q\),
and an integer target \(\ell\). It asks for sets \(X'\subseteq X\) and
\(Y'\subseteq Y\), with \(|X'|\le a\) and \(|Y'|\le q\), such that at least
\(\ell\) edges have an endpoint in \(X'\cup Y'\).
Write \(n_0=|X|\), \(m=|Y|\), and \(e=|E_H|\).
The source is unweighted; the two budgets are separate cardinality caps.

\paragraph{The regular-Clique source map.}
\label{hard:clique-source}
The following paragraphs detail the preceding reduction in the cited
article and its boundary cases. Let \(G\) be a simple \(d\)-regular graph with \(N\)
vertices, \(m_G\) edges, and nonnegative Clique target \(t\). Cases
\(t<3\), \(t>N\), and \(t=N\) are decidable in polynomial time; the
last case tests completeness. Cases \(d<t-1\) or
\(m_G<\binom{t}{2}\) are no-instances. Map every directly decided case
to the three-event temporal yes/no template below, carrying the numerical
budget \(b=t^2\). It remains to consider
\(3\le t<N\), \(d\ge t-1\), and \(m_G\ge\binom{t}{2}\).

Make one vertex of \(X\) for each vertex of \(G\), and one vertex of
\(Y\) for each edge of \(G\). Join \(v\in X\) to \(g\in Y\) exactly
when \(v\) is not an endpoint of \(g\). This non-incidence graph has
\(N+m_G\) vertices and \(e=(N-2)m_G\) edges. Set
\[
 a=t,\qquad q=\binom{t}{2},\qquad
 \ell=t(m_G-d)+(N-t)q.
\]
Any source selection within these caps can be extended to exactly
\(t\) left vertices and \(q\) right vertices without losing covered
edges. For these saturated choices \(S\subseteq V(G)\) and
\(J\subseteq E(G)\), let
\(i(S,J)=\sum_{g\in J}|\operatorname{endpoints}(g)\cap S|\).
Inclusion--exclusion gives the exact covered-edge count
\begin{equation}
\label{hard:clique-coverage}
\begin{aligned}
 C(S,J)
 &=t(m_G-d)+q(N-2)-\bigl(tq-i(S,J)\bigr)\\
 &=\ell-\bigl(2q-i(S,J)\bigr).
\end{aligned}
\end{equation}
Since \(i(S,J)\le2q\), reaching \(\ell\) is equivalent to every edge
of \(J\) having both endpoints in \(S\). A simple graph on \(t\)
vertices has at most \(q=\binom{t}{2}\) edges, so every such \(S\) is a
clique. Conversely, a \(t\)-clique and its internal
edges attain equality. This is the reduction of
\cite[Theorem~10]{mkrtchyan2024}, now with its at-most-to-exact step and
coverage identity explicit. The construction is polynomial, both side
budgets depend only on \(t\), and the temporal map below gives
\(b=a+2q=t^2\). If the resulting partial-cover target exceeds \(e\),
the following preprocessing correctly returns a no-instance.

\paragraph{Trivial partial-cover instances.}
The cases \(\ell\le0\) and \(\ell>e\) are immediate. If \(a\ge n_0\) or
\(q\ge m\), selecting an entire side covers every edge. If \(q=0\), choose
the \(a\) largest left degrees, or all left vertices when fewer are
available: edges incident to different left vertices are distinct, so this
maximizes coverage. Thus the construction may assume
\begin{equation}
\label{hard:source-range}
 0\le a\le n_0,\qquad 1\le q<m,\qquad 0\le\ell\le e.
\end{equation}
To preserve the exact formula \(b=a+2q\) on the directly decided cases as
well, use the three-event path
\(s\mathbin{\to}v\mathbin{\to}u\mathbin{\to}z\), at times \(0,1,2\),
with the same numerical budget \(b=a+2q\). Set \((T_2,T_3)=(1,1)\)
for a yes-instance and \((T_2,T_3)=(2,1)\) for a no-instance.
The first accepts \(D=\varnothing\); the second is impossible because its
undeleted two-walk count is one. Every event in either instance is
optional.

\paragraph{Rows, numerical constants, and columns.}
Add \(h=2m+1\) isolated left vertices to \(H\), and number the resulting
rows \(1,\ldots,n\). Define
\begin{equation}
\label{hard:constants}
\begin{aligned}
 n&=n_0+2m+1, & r&=n-a, & b&=a+2q,\\
 \kappa&=e+1, & P&=4m^2+1,\\
 \Lambda&=\bigl((\kappa+1)+r\kappa\bigr)mP+e+1 .
\end{aligned}
\end{equation}
In particular,
\begin{equation}
\label{hard:gaps}
\begin{gathered}
 r>2m,\qquad \kappa>e,\qquad P>2m^2,\\
 \Lambda+1>(\kappa+1)mP,\qquad
 \Lambda>r\kappa mP+e.
\end{gathered}
\end{equation}
The isolated rows have no source edges and require no additional budget.

For each \(y_j\in Y\), introduce an original column of numerical weight
\(j\) and a companion column of numerical weight \(P-j\).
All \(2m\) weights are distinct: the original weights are \(1,\ldots,m\),
and \(P-m>m\). Sort the columns by increasing weight and write their weights
as \(\rho_1<\cdots<\rho_{2m}\). Their sum is \(mP\).
These weights specify event counts in the construction; they are not event
deletion costs.

For row \(i\) and column \(t\), let \(N_{it}\) be the adjacency indicator
of the padded graph when column \(t\) is original. Set \(N_{it}=0\) for
companion columns and for padded rows. Let
\(\sigma_t=\sum_{i=1}^n N_{it}\); hence \(0\le\sigma_t\le n_0\).

\paragraph{The intermediate temporal instance.}
Use the static arcs
\[
 s_i\longrightarrow v_i\longrightarrow u\longrightarrow z
 \quad (1\le i\le n),\qquad v_0\longrightarrow u.
\]
The root is \(z\). There is a single depth-one vertex \(u\); the auxiliary
depth-two vertex \(v_0\) has no predecessor. For each row \(i\), put one
event \(A_i:s_i\to v_i\) at time zero. For every sorted column \(t\), put one
event \(C_t:u\to z\) at time \(2t\).

Prescribe the following cumulative numbers of middle events strictly
before \(C_t\):
\begin{equation}
\label{hard:prefixes}
 \mu_{it}=\Lambda+\kappa\rho_t+N_{it}
 \quad(1\le i\le n),\qquad
 \mu_{0t}=n+n\rho_t-\sigma_t .
\end{equation}
Set \(\mu_{i0}=0\) for \(0\le i\le n\).
For each \(t\), put exactly \(\mu_{it}-\mu_{i,t-1}\) distinct events
on \(v_i\to u\), all at time \(2t-1\), for every \(0\le i\le n\).
Call these the \(B\)-events.

All these multiplicities are nonnegative integers. The initial prefixes
are positive. For consecutive columns after the first, their increments
are bounded below by
\[
 \kappa(\rho_t-\rho_{t-1})+(N_{it}-N_{i,t-1})
 \ge\kappa-1\ge0
\]
on a productive row, and by
\[
 n(\rho_t-\rho_{t-1})-(\sigma_t-\sigma_{t-1})
 \ge n-n_0>0
\]
on the auxiliary row. The times \(2t-1\) and \(2t\) realize precisely the
prefixes in \eqref{hard:prefixes}, with strict temporal inequalities.
All contacts, including contacts sharing endpoints and a timestamp, receive
distinct IDs.

For the moment, \(B\)-events must survive, while \(A\)- and \(C\)-events are
optional. This is an intermediate construction only; the final instance
will make every event optional.

\begin{lemma}[Exact counts in the intermediate instance]
\label{hard:counts}
Let \(I\subseteq\{1,\ldots,n\}\) be the retained \(A\)-rows and
\(J\subseteq\{1,\ldots,2m\}\) the retained \(C\)-columns. Put
\[
 p=|I|,\qquad s=|J|,\qquad
 \omega=\sum_{t\in J}\rho_t,\qquad
 N(I,J)=\sum_{\substack{i\in I\\t\in J}}N_{it}.
\]
When all \(B\)-events survive, the counts and deletion cost are
\begin{equation}
\label{hard:identities}
\begin{aligned}
 W_2^\circ(I,J)
   &=n\bigl[s(\Lambda+1)+(\kappa+1)\omega\bigr],\\
 W_3^\circ(I,J)
   &=p\bigl[s\Lambda+\kappa\omega\bigr]+N(I,J),\\
 |D_{AC}|&=(n-p)+(2m-s).
\end{aligned}
\end{equation}
\end{lemma}

\begin{proof}
Each two-walk ending at \(z\) consists of a \(B\)-event followed by a
retained \(C_t\). Its contribution at column \(t\) is
\[
 \sum_{i=1}^n\mu_{it}+\mu_{0t}
 =n\Lambda+n\kappa\rho_t+\sigma_t
    +n+n\rho_t-\sigma_t
 =n\bigl[(\Lambda+1)+(\kappa+1)\rho_t\bigr].
\]
An \(A\)-event cannot belong to a two-walk ending at the root.
Each three-walk consists of \(A_i,B,C_t\), with \(i\in I\), \(t\in J\),
and the \(B\)-event counted by \(\mu_{it}\). The unique \(A_i\) precedes
every \(B\)-event on its row. The auxiliary row contributes no three-walks.
Summing \(\mu_{it}\) over \(I\times J\) proves the second identity. The last
identity counts the deleted optional IDs.
\end{proof}

Define the target retained-column count and target weight by
\[
 s_\star=2(m-q),\qquad \omega_\star=(m-q)P.
\]
They satisfy \(s_\star\ge2\), \(\omega_\star>0\), and
\(n+2m-b=r+s_\star\). Set the intermediate quotas to
\begin{equation}
\label{hard:intermediate-quotas}
\begin{aligned}
 T^\circ
   &=n\bigl[s_\star(\Lambda+1)+(\kappa+1)\omega_\star\bigr],\\
 U^\circ
   &=r\bigl[s_\star\Lambda+\kappa\omega_\star\bigr]+e-\ell.
\end{aligned}
\end{equation}

\begin{lemma}[The two quotas enforce the budget allocation]
\label{hard:allocation}
An \(A/C\) deletion set of cost at most \(b\) satisfies
\(W_2^\circ\ge T^\circ\) and \(W_3^\circ\le U^\circ\) if and only if
\begin{equation}
\label{hard:forced-values}
 p=r,\qquad s=s_\star,\qquad \omega=\omega_\star,\qquad
 N(I,J)\le e-\ell.
\end{equation}
Every feasible intermediate deletion set has cost exactly \(b\).
\end{lemma}

\begin{proof}
Suppose the two quotas hold. Since \(0\le\omega\le mP\), if
\(s<s_\star\) then
\[
 \frac{W_2^\circ-T^\circ}{n}
 =(s-s_\star)(\Lambda+1)
     +(\kappa+1)(\omega-\omega_\star)
 \le-(\Lambda+1)+(\kappa+1)mP<0,
\]
contradicting the lower quota. Therefore \(s\ge s_\star\).

Suppose instead that \(s=s_\star+d\), where \(d\ge1\).
The deletion budget gives \(p+s\ge r+s_\star\), so \(p\ge r-d\).
Since \(d\le s\le2m<r\), the latter lower bound is positive. Moreover,
\[
 ps\ge(r-d)(s_\star+d)
    =rs_\star+d(r-s)\ge rs_\star+1.
\]
All terms in the formula for \(W_3^\circ\) are nonnegative. It follows that
\[
\begin{aligned}
 W_3^\circ
 &\ge (rs_\star+1)\Lambda\\
 &>rs_\star\Lambda+r\kappa mP+e\\
 &\ge rs_\star\Lambda+r\kappa\omega_\star+e-\ell
   =U^\circ,
\end{aligned}
\]
a contradiction. Consequently \(s=s_\star\).

The budget now gives \(p\ge r\), and the lower quota gives
\(\omega\ge\omega_\star\). Subtracting the base term of \(U^\circ\)
from the triple count yields
\[
\begin{aligned}
 W_3^\circ-r(s_\star\Lambda+\kappa\omega_\star)
  ={}&(p-r)(s_\star\Lambda+\kappa\omega_\star)\\
    &+p\kappa(\omega-\omega_\star)+N(I,J).
\end{aligned}
\]
All three terms on the right are nonnegative, and their sum is at most
\(e-\ell\le e\). If \(p>r\), the first term is at least
\(s_\star\Lambda\ge2\Lambda>e\). If \(\omega>\omega_\star\), integrality
makes the second term at least \(r\kappa>e\).
Thus \(p=r\), \(\omega=\omega_\star\), and the remaining inequality is
\(N(I,J)\le e-\ell\). Conversely, substituting
\eqref{hard:forced-values} into \eqref{hard:identities} verifies both quotas.
The same substitution gives cost \((n-r)+(2m-s_\star)=a+2q=b\).
\end{proof}

\begin{lemma}[Companion balance]
\label{hard:balance}
If \(s=s_\star\) and \(\omega=\omega_\star\), exactly \(m-q\) original
columns and \(m-q\) companion columns survive.
\end{lemma}

\begin{proof}
Let \(c\) be the number of retained companion columns. Their weights and
the original-column weights give
\[
 \omega=cP+\delta,\qquad
 \delta=\sum_{\substack{j:\ \text{original }j\\\text{survives}}}j
       -\sum_{\substack{j:\ \text{companion }j\\\text{survives}}}j.
\]
There are at most \(2m\) summands, each of absolute value at most \(m\),
so \(|\delta|\le2m^2<P\). As \(\omega=(m-q)P\), the integer
\((m-q-c)P=\delta\) can only be zero. Hence \(c=m-q\);
the remaining \(s_\star-c=m-q\) retained columns are original.
The argument does not require matching identities for retained originals
and companions.
\end{proof}

\begin{lemma}[Equivalence with retained middle events]
\label{hard:retained-equivalence}
The source partial-cover instance is feasible if and only if the
intermediate temporal instance admits a deletion set of cost at most
\(b\), with \(W_2^\circ\ge T^\circ\) and \(W_3^\circ\le U^\circ\).
\end{lemma}

\begin{proof}
First take a source cover \(X',Y'\). Extend its selected left vertices to
exactly \(a\) rows of the padded graph, and extend \(Y'\) to exactly \(q\)
right vertices. Both extensions are possible and preserve coverage.
Delete the corresponding \(a\) events \(A_i\). For every selected right
vertex \(y_j\), delete its original \(C\)-event and its matching companion
\(C\)-event. This uses exactly \(a+2q\) deletions.
The retained columns consist of \(m-q\) pairs, each of weight \(P\).
Thus \(p=r\), \(s=s_\star\), and \(\omega=\omega_\star\).
The quantity \(N(I,J)\) is exactly the number of source edges not covered
by the extended selections, so it is at most \(e-\ell\).
Lemma~\ref{hard:allocation} gives the two temporal quotas.

Conversely, consider a feasible intermediate deletion set.
Lemma~\ref{hard:allocation} gives exactly \(a\) deleted \(A\)-events, and
Lemma~\ref{hard:balance} gives exactly \(q\) deleted original \(C\)-events
and \(q\) deleted companion \(C\)-events. Select the original left vertices
whose \(A\)-events were deleted, omitting any deleted padded rows, and
select the right vertices whose original \(C\)-events were deleted.
These selections use at most \(a\) left vertices and exactly \(q\) right
vertices. Only original columns and original rows can have nonzero
\(N_{it}\). Therefore \(N(I,J)\) counts precisely the source edges with
neither endpoint selected. Since \(N(I,J)\le e-\ell\), the selections
cover at least \(\ell\) edges.
\end{proof}

\paragraph{Removing every retention requirement.}
Set
\begin{equation}
\label{hard:replication-parameters}
 Q=2mb+1,\qquad T_2=QT^\circ,\qquad T_3=QU^\circ.
\end{equation}
Replace each intermediate \(B\)-event by \(Q\) distinct optional events
with the same endpoints and timestamp. Keep the \(A\)- and \(C\)-events
unchanged, and retain budget \(b\). All final events are original input
IDs with unit deletion cost. The intermediate instance is used only to
prove this final construction; it imposes no retention restriction on the
final input.

\begin{lemma}[Optional replication]
\label{hard:replication}
The all-optional final instance is feasible if and only if the
intermediate instance is feasible. In fact, every feasible final deletion
set deletes no \(B\)-event.
\end{lemma}

\begin{proof}
Fix the \(A/C\) choices of an arbitrary final deletion set \(D\) with
\(|D|\le b\), and let \(w_j\) be the intermediate count with those
choices and all middle events retained. Restoring every deleted \(B\)
copy in the final instance gives exactly \(Qw_j\) surviving \(j\)-walks,
for \(j=2,3\).

A \(B\) copy belongs to at most \(2m\) surviving three-walks: its row has
at most one \(A\)-event, and there are at most \(2m\) final \(C\)-events.
An auxiliary copy belongs to none. Each triple contains exactly one
\(B\) copy, so the triple loss caused by all deleted \(B\) copies is at
most
\[
 2m\,|D\cap B|\le2mb=Q-1.
\]
If \(w_3>U^\circ\), integrality would therefore imply
\[
 W_3(E\setminus D)\ge Q(U^\circ+1)-(Q-1)=QU^\circ+1>T_3,
\]
contrary to feasibility. Hence \(w_3\le U^\circ\).
Restoring middle events can only increase the two-walk count; consequently
\[
 Qw_2\ge W_2(E\setminus D)\ge QT^\circ,
\]
and \(w_2\ge T^\circ\).
The \(A/C\) part of \(D\) alone costs at most \(b\), so it is a feasible
intermediate solution. By Lemma~\ref{hard:allocation}, that part already
uses exactly \(b\) deletions. Thus \(D\) contains no deleted \(B\) copy.

Conversely, a feasible intermediate \(A/C\) deletion set can be used in
the final instance while deleting no \(B\) copy. The two counts are then
exactly \(Qw_2\) and \(Qw_3\), and both final quotas hold.
\end{proof}

\begin{lemma}[Encoded size and parameter map]
\label{hard:size}
The nontrivial construction contains exactly
\begin{equation}
\label{hard:event-count}
 M=n+2m+Qn\bigl[\Lambda+(\kappa+1)P-\kappa\bigr]
\end{equation}
explicitly listed events, uses timestamps in \(\{0,\ldots,4m\}\), and
has size polynomial in the explicit source input length.
Its deletion budget is \(b=a+2q\), and the composition with the
\(t\)-Clique reduction in \cite[Theorem~10]{mkrtchyan2024} has \(b=t^2\).
\end{lemma}

\begin{proof}
The last column has weight \(P-1\) and is a companion; hence all its
adjacency indicators and its \(\sigma\)-value are zero. Each productive
row contains \(\Lambda+\kappa(P-1)\) intermediate \(B\)-events, and the
auxiliary row contains \(nP\). Their total is
\[
 n\bigl[\Lambda+\kappa(P-1)\bigr]+nP
   =n\bigl[\Lambda+(\kappa+1)P-\kappa\bigr].
\]
Multiplying by \(Q\) and adding \(n\) events of type \(A\) and \(2m\)
of type \(C\) gives \eqref{hard:event-count}.
The static tree has \(2n+3\) vertices and depth three. Its only
depth-one vertex is \(u\). Replication changes neither the static tree
nor any timestamp.

Put \(S_H=n_0+m\). In the range \eqref{hard:source-range},
\[
\begin{gathered}
 e=O(S_H^2),\quad n,r,b=O(S_H),\quad
 \kappa,P,Q=O(S_H^2),\\
 \Lambda=O(S_H^6),\quad M=O(S_H^9).
\end{gathered}
\]
Moreover \(T^\circ,U^\circ=O(S_H^8)\), so \(T_2,T_3=O(S_H^{10})\).
Event IDs, endpoints, timestamps, budgets, and quotas therefore need
only \(O(\log S_H)\) bits per numerical field. In particular, an explicit
event list has length \(O(S_H^9\log S_H)\). Every contact is written out;
the construction uses no binary-encoded multiplicities.
The directly decided cases have constant-size temporal event lists and
retain the binary encoding of \(b=a+2q\).

The parameter bound \(b=a+2q\le2(a+q)\) is independent of \(n_0,m,e\).
The source theorem uses budgets
\(a=t\) and \(q=\binom{t}{2}\) in its reduction from regular
\(t\)-Clique. Thus the composed budget is
\[
 b=t+2\binom{t}{2}=t^2.
\]
This identity also holds if a directly decided source case uses the
three-event template described above.
\end{proof}

\paragraph{Why the preservation-slack algorithm is consistent with this result.}
For the nontrivial construction, the final instance has no mandatory
deletions and \(T_2>0\).
Its preservation slack \(R=W_2(E)-T_2\) is exactly
\begin{equation}
\label{hard:slack}
 R=Qnq\bigl[2(\Lambda+1)+(\kappa+1)P\bigr].
\end{equation}
Indeed, the full retained column set has \(s=2m\) and weight \(mP\);
subtracting \(QT^\circ\) from the first identity in
\eqref{hard:identities}, scaled by \(Q\), gives \eqref{hard:slack}.
This quantity is unbounded in the source graph size even with fixed
budgets. For example, fix \(a=q=1\), \(m=2\), \(e=\ell=1\), and increase
the number of isolated left vertices. Then \(b=3\) is fixed, whereas
\(n\), \(\Lambda\), and \(R\) grow.
Thus this reduction does not bound \(R\) by a function of \(b\).
It is consistent with a \(2^{O(\sqrt R)}\operatorname{poly}(|\mathcal I|)\)
algorithm parameterized by preservation slack. The conclusion is
\(\mathrm{W}[1]\)-hardness in \(b\); no \(\mathrm{W}[1]\)-membership claim
or ETH exponent lower bound is asserted.
% END COMPLETE CONSTRUCTION.

\subsection{Classical complexity consequence}

The explicit construction also yields the following classical consequence.
The proof is included for completeness, without a historical-priority claim
for this corollary or the elementary regularization used below.

\begin{corollary}[Classical NP-completeness]
\label{full:np-complete}
The decision problem in Theorem~\ref{hard:budget} is NP-complete, even with
all events optional and of unit cost, one depth-one vertex, and exactly
one time-zero \(A\) event at each productive depth-two vertex.
\end{corollary}

\begin{proof}
Membership in NP follows from an original-event deletion bit vector.
Its budget and forced statuses, when present, can be checked directly.
Enumerating ordered retained event pairs and triples, and testing their
endpoints and strict timestamp inequalities, computes the required counts
using \(O(M^3)\) candidate tuples. Endpoint and timestamp comparisons,
comparison with the encoded quotas, and validation of the structural
restrictions all take polynomial time in the complete input length.

Classical hardness follows by a polynomial chain of reductions from ordinary Clique.
Let \(G\) have \(N\) vertices and maximum degree \(\Delta\), with target
\(t\). Cases \(t<3\), \(t>N\), or \(\Delta=0\) are directly decidable.
Otherwise make \(2\Delta\) disjoint copies of \(G\), named \(L_i,R_i\)
for \(i\in\mathbb Z_\Delta\). For every original vertex \(v\) of degree
\(d(v)\), join its copy in \(L_i\) to its copies in \(R_{i+j}\), for
\(0\le j<\Delta-d(v)\), with indices modulo \(\Delta\).
Each copied vertex receives exactly its degree deficit, on both sides,
so the resulting simple graph is \(\Delta\)-regular. No added edge lies
in a triangle: its endpoints have the same original label in different
copies; a vertex of another label can be an internal neighbor of at most
one endpoint, while all edges among copies of the same label are
bipartite. Thus every clique of size at least three remains within an
original copy, preserving the existence of a \(t\)-clique.
The regularization has \(2N\Delta=O(N^2)\) vertices and
\(N\Delta^2=O(N^3)\) edges and is polynomial.

The regularized graph has \(N'=2N\Delta>t\), so the strict source-size
condition is automatic. Apply the regular-Clique prechecks and source-map
argument above. For its \(N'\) vertices, \(m'\) edges, and degree \(d\),
the non-incidence construction in \cite[Theorem~10]{mkrtchyan2024}
uses one source-side vertex per vertex and one per edge, with
\[
 a=t,\qquad q=\binom{t}{2},\qquad
 \ell=(N'-t)\binom{t}{2}+t(m'-d).
\]
The graph has \(N'+m'\) vertices and at most \(N'm'\) edges, so this
step is polynomial. Apply the temporal construction above.
Its side-size sum satisfies \(S_H=O(N^3)\); hence
Lemma~\ref{hard:size} gives polynomial encoded output size, even with
the loose event bound \(O(N^{27})\). Directly decided source cases use
the stated constant-event yes/no templates. The output has the promised
structural restrictions, and its budget is
\(b=t+2\binom{t}{2}=t^2\).
This polynomial many-one reduction proves NP-hardness and completes
the proof.
\end{proof}

Consequently, an exact polynomial-time algorithm for every instance of
this full model would imply \(\mathrm P=\mathrm{NP}\). This consequence
does not rule out a simpler or faster algorithm parameterized by
preservation slack \(R\), and it supplies no matching lower bound for
the square-root exponent.
